\documentclass[11pt,a4paper]{article}

\usepackage[utf8]{inputenc}
\usepackage[T1]{fontenc}
\usepackage{geometry}
\usepackage{hyperref}
\usepackage{longtable}
\usepackage{booktabs}
\usepackage{array}
\usepackage{listings}
\usepackage{xcolor}
\usepackage{enumitem}

\geometry{margin=2.2cm}
\hypersetup{
    colorlinks=true,
    linkcolor=blue!50!black,
    urlcolor=blue!50!black
}

\lstset{
    basicstyle=\ttfamily\small,
    breaklines=true,
    columns=fullflexible,
    keepspaces=true,
    showstringspaces=false,
    frame=single,
    rulecolor=\color{black!20}
}

\title{Evalon Backtest API\\Kullanim Dokumani}
\author{Codex tarafindan kod tabani uzerinden uretilmistir}
\date{17 Nisan 2026}

\begin{document}
\maketitle
\tableofcontents
\newpage

\section{Amac ve Kapsam}

Bu dokuman, \texttt{backend/backtest} klasorundeki FastAPI tabanli servisin tum ana ozelliklerini, endpointlerini, konfigrasyonunu ve kullanim seklini aciklar. Servis esas olarak dort ana yetenek sunar:

\begin{itemize}[leftmargin=1.5em]
    \item BIST/Oracle tabanli fiyat verisi sunma
    \item TA-Lib uzerinden indikator hesaplama
    \item Blueprint tabanli backtest calistirma
    \item AI yardimli strateji, kural ve indikator taslagi uretme
\end{itemize}

Bu dokuman kod tabanindaki mevcut uygulama davranisina gore yazilmistir. Ozellikle \texttt{api/main.py}, \texttt{api/modules/backtests/*} ve \texttt{api/modules/ai/*} altindaki kaynaklar esas alinmistir.

\section{Genel Mimari}

\subsection{Servisin Rolü}

Bu API tek bir FastAPI uygulamasi olarak calisir ve asagidaki fonksiyonlari tek cati altinda sunar:

\begin{itemize}[leftmargin=1.5em]
    \item \texttt{/v1/prices}: mum verisi
    \item \texttt{/v1/indicators*}: indikator katalogu ve indikator hesaplama
    \item \texttt{/v1/backtests*}: rule katalogu, preset katalogu ve backtest calistirma
    \item \texttt{/v1/ai*}: AI session, mesajlasma, draft ve asset kayitlari
\end{itemize}

\subsection{Veri Kaynagi}

Fiyat verisi Oracle uzerindeki BIST verisinden okunur. Uygulama \texttt{HybridBistPricesClient} kullanir:

\begin{itemize}[leftmargin=1.5em]
    \item dusuk timeframe'lerde klasik tablo sorgulari kullanilir
    \item \texttt{1h} ve ustunde saatlik veri tabani tablosu tercih edilir
    \item Oracle wallet desteklenir
    \item connection pool opsiyonel olarak acilir
\end{itemize}

\subsection{Test Modu}

\texttt{ticker=TEST} ozel bir davranistir. Bu durumda Oracle yerine repo icindeki fixture dosyasi okunur:

\begin{itemize}[leftmargin=1.5em]
    \item varsayilan yol: \texttt{api/fixtures/test\_ohlcv.csv}
    \item amac: Oracle credentials olmadan smoke test ve entegrasyon testi yapabilmek
\end{itemize}

\subsection{State ve Kalicilik}

v1 surumunde tum state kalici degildir:

\begin{itemize}[leftmargin=1.5em]
    \item backtest run bilgileri: \texttt{InMemoryRunStore}
    \item AI session'lari: \texttt{InMemoryAiSessionStore}
    \item AI asset kayitlari: JSON dosyasi veya gerekli durumda memory fallback
\end{itemize}

Bu nedenle Cloud Run instance restart, revision degisimi veya coklu instance senaryolarinda run/session kayitlari kaybolabilir. Bu davranis ozellikle \texttt{/v1/backtests/start} ve \texttt{/v1/ai/sessions/*} kullaniminda dikkate alinmalidir.

\section{Calisma ve Konfigurasyon}

\subsection{Calistirma}

Lokal gelistirme:

\begin{lstlisting}[language=bash]
python3 -m venv venv
./venv/bin/pip install -r requirements.txt
./venv/bin/uvicorn api.main:app --reload --host 0.0.0.0 --port 8000
\end{lstlisting}

Cloud Run icin Docker tabanli deploy dosyalari mevcuttur:

\begin{itemize}[leftmargin=1.5em]
    \item \texttt{Dockerfile}
    \item \texttt{scripts/deploy\_cloud\_run.sh}
    \item \texttt{scripts/smoke\_test\_cloud\_run.sh}
    \item \texttt{DEPLOY\_CLOUD\_RUN.md}
\end{itemize}

Bu dokuman hazirlandiginda aktif Cloud Run deployment ornegi:

\begin{quote}
\url{https://evalon-backtest-api-474112640179.europe-west1.run.app}
\end{quote}

\subsection{Zorunlu ve Opsiyonel Environment Variable'lar}

\begin{longtable}{>{\raggedright\arraybackslash}p{4.5cm} >{\raggedright\arraybackslash}p{2.4cm} >{\raggedright\arraybackslash}p{7.5cm}}
\toprule
Degisken & Durum & Aciklama \\
\midrule
\endhead
\texttt{ORACLE\_DB\_USER} & Zorunlu & Oracle kullanici adi. Ornek: \texttt{ADMIN}. \\
\texttt{ORACLE\_DB\_PASSWORD} & Zorunlu & Oracle baglanti sifresi. \\
\texttt{ORACLE\_DB\_DSN} & Zorunlu & Oracle TNS alias'i. Ornek: \texttt{evalondb\_high}. \\
\texttt{ORACLE\_WALLET\_ZIP\_B64} & Cloud icin onerilen & Base64 zip formatinda Oracle wallet. Cloud Run/Vercel gibi ortamlarda en pratik yontem budur. \\
\texttt{ORACLE\_WALLET\_DIR} & Opsiyonel & Wallet klasor yolu. Lokal calismada kullanisli. \\
\texttt{ORACLE\_WALLET\_PASSWORD} & Opsiyonel & Wallet sifresi gerekiyorsa kullanilir. \\
\texttt{ORACLE\_USE\_POOL} & Opsiyonel & \texttt{1} ise Oracle connection pool acilir. Varsayilan \texttt{1}. \\
\texttt{ORACLE\_POOL\_MIN} & Opsiyonel & Pool min baglanti sayisi. \\
\texttt{ORACLE\_POOL\_MAX} & Opsiyonel & Pool max baglanti sayisi. \\
\texttt{ORACLE\_POOL\_INC} & Opsiyonel & Pool artim degeri. \\
\texttt{BIST\_DEBUG} & Opsiyonel & \texttt{1} ise detayli hata detaylari donebilir. \\
\texttt{AI\_ASSET\_STORE\_PATH} & Opsiyonel & AI asset JSON dosyasinin yolu. Cloud Run tarafinda tipik deger \texttt{/tmp/evalon\_ai\_assets.json}. \\
\texttt{GEMINI\_API\_KEY} & Opsiyonel & AI modulu icin Gemini API key. Yoksa heuristic fallback veya Vertex kullanimi denenir. \\
\texttt{GOOGLE\_CLOUD\_PROJECT} & Opsiyonel & Vertex/GCP baglami icin proje adi. \\
\texttt{GOOGLE\_CLOUD\_LOCATION} & Opsiyonel & Vertex lokasyonu. Varsayilan \texttt{global}. \\
\texttt{ALLOWED\_ORIGINS} & Opsiyonel & Virgulle ayrilmis CORS origin listesi. \\
\texttt{ALLOWED\_ORIGIN\_REGEX} & Opsiyonel & Regex tabanli CORS ayari. \\
\texttt{TEST\_OHLCV\_PATH} & Opsiyonel & \texttt{ticker=TEST} senaryosunda okunacak CSV veya ZST dosyasi. \\
\bottomrule
\end{longtable}

\subsection{CORS Davranisi}

Uygulama uretim ortami icin wildcard CORS kullanmaz. Beklenen model:

\begin{itemize}[leftmargin=1.5em]
    \item \texttt{ALLOWED\_ORIGINS} ver
    \item veya \texttt{ALLOWED\_ORIGIN\_REGEX} ver
\end{itemize}

Lokal ortamda explicit CORS verilmemisse varsayilan olarak \texttt{localhost:3000}, \texttt{127.0.0.1:3000}, \texttt{localhost:3001}, \texttt{127.0.0.1:3001} izinli kabul edilir.

\section{Genel API Kurallari}

\subsection{Auth}

Mevcut endpointlerin hicbirinde HTTP auth zorunlulugu yoktur. Bu nedenle servis su an public olarak calisir. Uretim guvenligi icin gerekirse reverse proxy veya uygulama seviyesi auth eklenmelidir.

\subsection{Formatlar}

\begin{itemize}[leftmargin=1.5em]
    \item tum endpointler JSON doner
    \item zaman alanlari genellikle ISO 8601 tarih-zaman veya epoch second olarak gelir
    \item hata durumlarinda FastAPI \texttt{detail} alani iceren standart JSON hata yapisi doner
\end{itemize}

\subsection{Yaygin HTTP Durum Kodlari}

\begin{itemize}[leftmargin=1.5em]
    \item \texttt{200}: basarili
    \item \texttt{400}: request body veya parametre dogrulama hatasi
    \item \texttt{404}: run ya da session bulunamadi
    \item \texttt{500}: Oracle, TA-Lib, dosya sistemi veya altyapi hatasi
\end{itemize}

\section{Endpoint Ozet Tablosu}

\begin{longtable}{>{\raggedright\arraybackslash}p{2.4cm} >{\raggedright\arraybackslash}p{5.8cm} >{\raggedright\arraybackslash}p{6cm}}
\toprule
Method & Path & Amac \\
\midrule
\endhead
\texttt{GET} & \texttt{/} & Kisa landing payload ve config snapshot \\
\texttt{GET} & \texttt{/health} & Health check \\
\texttt{GET} & \texttt{/v1/prices} & Fiyat/mum verisi \\
\texttt{GET} & \texttt{/v1/indicators/catalog} & Desteklenen indikatorler \\
\texttt{GET} & \texttt{/v1/indicators} & Indikator serisi hesaplama \\
\texttt{GET} & \texttt{/v1/backtests/catalog/rules} & Rule katalogu \\
\texttt{GET} & \texttt{/v1/backtests/catalog/presets} & Preset katalogu \\
\texttt{POST} & \texttt{/v1/backtests/run} & Senkron backtest \\
\texttt{POST} & \texttt{/v1/backtests/start} & Asenkron backtest \\
\texttt{GET} & \texttt{/v1/backtests/\{run\_id\}/status} & Run durumu \\
\texttt{GET} & \texttt{/v1/backtests/\{run\_id\}/events} & Event sayfalama \\
\texttt{GET} & \texttt{/v1/backtests/\{run\_id\}/portfolio-curve} & Portfoy egri verisi \\
\texttt{GET} & \texttt{/v1/ai/tools} & AI tool katalogu \\
\texttt{POST} & \texttt{/v1/ai/sessions} & Yeni AI session \\
\texttt{GET} & \texttt{/v1/ai/sessions/\{session\_id\}} & Session detaylari \\
\texttt{POST} & \texttt{/v1/ai/sessions/\{session\_id\}/messages} & AI ile mesajlasma \\
\texttt{GET} & \texttt{/v1/ai/assets} & Kayitli AI asset'leri \\
\texttt{POST} & \texttt{/v1/ai/strategies} & AI strategy asset kaydi \\
\texttt{POST} & \texttt{/v1/ai/rules} & AI rule asset kaydi \\
\texttt{POST} & \texttt{/v1/ai/indicators} & AI indicator asset kaydi \\
\bottomrule
\end{longtable}

\section{Temel Endpointler}

\subsection{\texttt{GET /}}

\paragraph{Ne ise yarar?}
Servisin root endpointidir. 404 yerine anlamli bir JSON doner ve servis yuzeyini gosterir.

\paragraph{Ne doner?}
\begin{itemize}[leftmargin=1.5em]
    \item servis adi
    \item health path'i
    \item docs path'i
    \item openapi path'i
    \item ornek prices path'i
    \item secret degerleri aciga cikarmayan config snapshot
\end{itemize}

\subsection{\texttt{GET /health}}

\paragraph{Ne ise yarar?}
Load balancer, uptime monitor ve Cloud Run smoke test icin kullanilir.

\paragraph{Ornek cevap}
\begin{lstlisting}
{
  "status": "ok"
}
\end{lstlisting}

\section{Fiyat ve Indikator Endpointleri}

\subsection{\texttt{GET /v1/prices}}

\paragraph{Amac}
Belirli bir ticker ve timeframe icin OHLCV mum verisi dondurur.

\paragraph{Query parametreleri}
\begin{longtable}{>{\raggedright\arraybackslash}p{3cm} >{\raggedright\arraybackslash}p{2.2cm} >{\raggedright\arraybackslash}p{8.2cm}}
\toprule
Parametre & Tip & Aciklama \\
\midrule
\endhead
\texttt{ticker} & string & Zorunlu. Ornek: \texttt{THYAO}. Ozel test modu icin \texttt{TEST}. \\
\texttt{timeframe} & string & Varsayilan \texttt{1m}. Ornekler: \texttt{1m}, \texttt{3m}, \texttt{5m}, \texttt{15m}, \texttt{30m}, \texttt{1h}, \texttt{2h}, \texttt{4h}, \texttt{6h}, \texttt{12h}, \texttt{1d}, \texttt{1g}, \texttt{1w}, \texttt{1M}, \texttt{1mo}. \\
\texttt{start} & datetime & Opsiyonel alt sinir. \\
\texttt{end} & datetime & Opsiyonel ust sinir. \\
\texttt{limit} & int & Opsiyonel. \texttt{1} ile \texttt{200000} arasi. \\
\bottomrule
\end{longtable}

\paragraph{Ornek cURL}
\begin{lstlisting}[language=bash]
curl "http://127.0.0.1:8000/v1/prices?ticker=THYAO&timeframe=1h&limit=24"
\end{lstlisting}

\paragraph{Ornek cevap}
\begin{lstlisting}
{
  "ticker": "THYAO",
  "timeframe": "1h",
  "rows": 24,
  "data": [
    {
      "t": "2026-04-17T09:00:00+00:00",
      "o": 287.4,
      "h": 289.1,
      "l": 286.8,
      "c": 288.7,
      "v": 345678
    }
  ]
}
\end{lstlisting}

\paragraph{Ozel davranislar}
\begin{itemize}[leftmargin=1.5em]
    \item \texttt{ticker=TEST} ise Oracle yerine fixture dosyasi okunur.
    \item Oracle sifresi yoksa \texttt{500} doner.
    \item Wallet configure edilmemisse \texttt{500} doner.
    \item gecersiz timeframe gibi durumlarda \texttt{400} doner.
\end{itemize}

\subsection{\texttt{GET /v1/indicators/catalog}}

\paragraph{Amac}
Desteklenen indikatorlerin katalogunu dondurur. UI tarafinda dropdown doldurmak veya hangi strategy slug'larin desteklendigini ogrenmek icin kullanilir.

\paragraph{Ornek cevap}
\begin{lstlisting}
{
  "count": 30,
  "indicators": [
    { "id": "rsi", "label": "RSI" },
    { "id": "macd", "label": "MACD" }
  ]
}
\end{lstlisting}

\subsection{\texttt{GET /v1/indicators}}

\paragraph{Amac}
Belirli bir ticker/timeframe uzerinde TA-Lib ile indikator serisi hesaplar.

\paragraph{Query parametreleri}
\begin{longtable}{>{\raggedright\arraybackslash}p{3cm} >{\raggedright\arraybackslash}p{2.2cm} >{\raggedright\arraybackslash}p{8.2cm}}
\toprule
Parametre & Tip & Aciklama \\
\midrule
\endhead
\texttt{ticker} & string & Zorunlu. \\
\texttt{timeframe} & string & Varsayilan \texttt{1m}. \\
\texttt{strategy} & string & Zorunlu indikator slug'i. Ornek: \texttt{rsi}, \texttt{macd}, \texttt{sma}, \texttt{ema}, \texttt{bbands}. \\
\texttt{period} & int & Ozellikle RSI ve moving average tabanli indikatorlerde kullanilir. Varsayilan \texttt{14}. \\
\texttt{fast} & int & MACD hizli period. Varsayilan \texttt{12}. \\
\texttt{slow} & int & MACD yavas period. Varsayilan \texttt{26}. \\
\texttt{signal} & int & MACD signal period. Varsayilan \texttt{9}. \\
\texttt{start} & datetime & Opsiyonel alt sinir. \\
\texttt{end} & datetime & Opsiyonel ust sinir. \\
\texttt{limit} & int & Opsiyonel. \texttt{1} ile \texttt{200000} arasi. \\
\bottomrule
\end{longtable}

\paragraph{Ornek cURL}
\begin{lstlisting}[language=bash]
curl "http://127.0.0.1:8000/v1/indicators?ticker=THYAO&timeframe=1h&strategy=rsi&period=14&limit=100"
\end{lstlisting}

\paragraph{Cevap yapisi}
\begin{lstlisting}
{
  "ticker": "THYAO",
  "timeframe": "1h",
  "strategy": "rsi",
  "indicators": [
    {
      "id": "rsi",
      "series": [
        { "t": "2026-04-17T09:00:00+00:00", "v": 58.2 }
      ]
    }
  ]
}
\end{lstlisting}

\paragraph{Notlar}
\begin{itemize}[leftmargin=1.5em]
    \item desteklenmeyen strategy slug'lari \texttt{400} ile reddedilir
    \item TA-Lib kurulu degilse \texttt{500} doner
    \item veri bos ise \texttt{indicators: []} doner
\end{itemize}

\section{Backtest Endpointleri}

\subsection{Genel Yapi}

Backtest modulu blueprint tabanlidir. Request body'si temel olarak su bloklardan olusur:

\begin{itemize}[leftmargin=1.5em]
    \item \texttt{symbol} veya \texttt{symbols}
    \item \texttt{stageThreshold}
    \item \texttt{direction}
    \item \texttt{testWindowDays}
    \item \texttt{portfolio}
    \item \texttt{risk}
    \item \texttt{stages.trend/setup/trigger}
\end{itemize}

Blueprint validasyon kurallari:

\begin{itemize}[leftmargin=1.5em]
    \item en az bir \texttt{symbol} veya \texttt{symbols} verilmeli
    \item \texttt{stageThreshold} 1 ile 3 arasinda olmali
    \item \texttt{stages} sozlugu zorunlu
    \item \texttt{trend}, \texttt{setup}, \texttt{trigger} stage'leri mevcut olmali
    \item secilen rule id'leri katalogda tanimli olmali
\end{itemize}

\subsection{\texttt{GET /v1/backtests/catalog/rules}}

\paragraph{Amac}
Rule katalogunu ve family bilgilerini dondurur. UI tarafinda rule secici olusturmak icin kullanilir.

\paragraph{Cevap ana alanlari}
\begin{itemize}[leftmargin=1.5em]
    \item \texttt{count}
    \item \texttt{families}
    \item \texttt{rules[]}: \texttt{id}, \texttt{label}, \texttt{family}, \texttt{category}, \texttt{stages}, \texttt{summary}
\end{itemize}

\subsection{\texttt{GET /v1/backtests/catalog/presets}}

\paragraph{Amac}
Hazir strateji preset'lerini dondurur.

\paragraph{Cevap ana alanlari}
\begin{itemize}[leftmargin=1.5em]
    \item \texttt{count}
    \item \texttt{presets[]}: \texttt{id}, \texttt{label}, \texttt{summary}, \texttt{direction}, \texttt{stageThreshold}, \texttt{ruleIds}
\end{itemize}

\subsection{\texttt{POST /v1/backtests/run}}

\paragraph{Amac}
Backtest'i senkron calistirir; request bitmeden sonuc dondurur.

\paragraph{Ne zaman kullanilir?}
Kisa ve aninda sonuc beklenen isler icin.

\paragraph{Ornek request body}
\begin{lstlisting}
{
  "symbol": "THYAO",
  "symbols": ["THYAO"],
  "stageThreshold": 2,
  "direction": "long",
  "testWindowDays": 365,
  "portfolio": {
    "initialCapital": 100000,
    "positionSize": 10000,
    "commissionPct": 0.1
  },
  "risk": {
    "stopPct": 1.8,
    "targetPct": 4.0,
    "maxBars": 12
  },
  "stages": {
    "trend": {
      "key": "trend",
      "timeframe": "1d",
      "required": false,
      "minOptionalMatches": 0,
      "rules": [
        {
          "id": "trend-slope",
          "required": true,
          "params": { "lookback": 20, "minMovePct": 2.0 }
        }
      ]
    },
    "setup": {
      "key": "setup",
      "timeframe": "4h",
      "required": true,
      "minOptionalMatches": 0,
      "rules": [
        {
          "id": "breakout",
          "required": true,
          "params": { "lookback": 20, "buffer": 0.2 }
        }
      ]
    },
    "trigger": {
      "key": "trigger",
      "timeframe": "1h",
      "required": true,
      "minOptionalMatches": 0,
      "rules": [
        {
          "id": "micro-breakout",
          "required": true,
          "params": { "bars": 4 }
        }
      ]
    }
  }
}
\end{lstlisting}

\paragraph{Cevap ana alanlari}
\begin{itemize}[leftmargin=1.5em]
    \item \texttt{runId}
    \item \texttt{result}
    \item \texttt{summary}: \texttt{totalTrades}, \texttt{winRate}, \texttt{totalPnl}, \texttt{maxDrawdown}
    \item \texttt{eventsCount}
\end{itemize}

\subsection{\texttt{POST /v1/backtests/start}}

\paragraph{Amac}
Backtest'i arka planda baslatir. Hemen \texttt{queued} doner; sonrasinda client status endpointleri ile ilerlemeyi izler.

\paragraph{Cevap ornegi}
\begin{lstlisting}
{
  "runId": "btrun_xxx",
  "status": "queued",
  "progress": {
    "phase": "queued",
    "progressPct": 0,
    "totalSymbols": 3,
    "processedSymbols": 0,
    "currentSymbol": null,
    "message": "Backtest kuyruga alindi."
  }
}
\end{lstlisting}

\paragraph{Onemli not}
Bu mekanizma arka planda Python thread kullanir ve run store memory tabanlidir. Process restart olursa run kaybolur.

\subsection{\texttt{GET /v1/backtests/\{run\_id\}/status}}

\paragraph{Amac}
Run'in mevcut durumunu verir.

\paragraph{Durum alanlari}
\begin{itemize}[leftmargin=1.5em]
    \item \texttt{queued}
    \item \texttt{running}
    \item \texttt{completed}
    \item \texttt{failed}
\end{itemize}

\paragraph{Cevap ana alanlari}
\begin{itemize}[leftmargin=1.5em]
    \item \texttt{runId}, \texttt{status}
    \item \texttt{createdAt}, \texttt{startedAt}, \texttt{finishedAt}
    \item \texttt{progress}
    \item \texttt{eventsCount}
    \item basariliysa \texttt{result} ve \texttt{summary}
    \item hata varsa \texttt{error}
\end{itemize}

\subsection{\texttt{GET /v1/backtests/\{run\_id\}/events}}

\paragraph{Amac}
Trade event'lerini sayfali dondurur.

\paragraph{Query parametreleri}
\begin{itemize}[leftmargin=1.5em]
    \item \texttt{page}: varsayilan 1, minimum 1
    \item \texttt{limit}: varsayilan 200, maksimum 500
\end{itemize}

\paragraph{Cevap ana alanlari}
\begin{itemize}[leftmargin=1.5em]
    \item \texttt{events[]}
    \item \texttt{summary}
    \item \texttt{page}, \texttt{totalPages}, \texttt{totalEvents}
\end{itemize}

Her event tipik olarak asagidaki alanlari tasir:

\begin{lstlisting}
{
  "id": "trade_1_entry",
  "type": "entry",
  "time": 1710000000,
  "price": 100.0,
  "side": "long",
  "symbol": "THYAO",
  "qty": 10.0,
  "tradeId": "trade_1",
  "orderId": "btrun_x_trade_1_entry",
  "pnl": null,
  "reason": "Blueprint trigger",
  "meta": {
    "allocatedCapital": 1000,
    "score": 3,
    "stages": { "trend": true, "setup": true, "trigger": true }
  }
}
\end{lstlisting}

Exit event tipi \texttt{tp}, \texttt{stop} veya \texttt{exit} olabilir.

\subsection{\texttt{GET /v1/backtests/\{run\_id\}/portfolio-curve}}

\paragraph{Amac}
Portfoy equity curve verisini dondurur.

\paragraph{Cevap ana alanlari}
\begin{itemize}[leftmargin=1.5em]
    \item status bilgisi
    \item progress bilgisi
    \item hata varsa \texttt{error}
    \item sonuc varsa \texttt{summary}
    \item \texttt{curve}: \texttt{portfolioCurve} verisi
\end{itemize}

\section{AI Endpointleri}

\subsection{AI Modulu Ne Yapar?}

AI modulu kullanicidan mesaj alip bunu bir execution plan'a donusturur, arka planda uygun tool'lari cagirir ve strateji/rule/indikator taslagi uretebilir.

AI tarafinda iki katman vardir:

\begin{itemize}[leftmargin=1.5em]
    \item \textbf{Session}: konusma gecmisi ve aktif context
    \item \textbf{Asset}: kullaniciya ait strategy/rule/indicator taslaklari
\end{itemize}

\subsection{AI Tool Katalogu}

\texttt{GET /v1/ai/tools} endpointi mevcut AI araclarini listeler. Kod tabanina gore temel tool seti su isimlerden olusur:

\begin{itemize}[leftmargin=1.5em]
    \item \texttt{get\_rule\_catalog}
    \item \texttt{get\_preset\_catalog}
    \item \texttt{get\_prices}
    \item \texttt{get\_indicator\_catalog}
    \item \texttt{get\_indicators}
    \item \texttt{run\_backtest}
    \item \texttt{get\_backtest\_status}
    \item \texttt{get\_backtest\_events}
    \item \texttt{get\_portfolio\_curve}
    \item \texttt{list\_user\_assets}
    \item \texttt{save\_user\_strategy}
    \item \texttt{save\_user\_rule}
    \item \texttt{save\_user\_indicator}
\end{itemize}

\subsection{\texttt{POST /v1/ai/sessions}}

\paragraph{Amac}
Yeni bir AI oturumu olusturur.

\paragraph{Request body}
\begin{lstlisting}
{
  "userId": "demo",
  "title": "Trend stratejisi arastirmasi"
}
\end{lstlisting}

\paragraph{Not}
Session store memory tabanlidir. Cloud Run restart olursa session kaybolabilir.

\subsection{\texttt{GET /v1/ai/sessions/\{session\_id\}}}

\paragraph{Amac}
Belirli bir AI session'in mevcut kaydini getirir.

\paragraph{Cevapta gorulebilecek alanlar}
\begin{itemize}[leftmargin=1.5em]
    \item \texttt{session\_id}
    \item \texttt{user\_id}
    \item \texttt{messages}
    \item \texttt{last\_plan}
    \item \texttt{last\_tool\_results}
    \item \texttt{working\_context}
    \item \texttt{last\_backtest\_run\_id}
    \item \texttt{active\_strategy\_draft}
\end{itemize}

\subsection{\texttt{POST /v1/ai/sessions/\{session\_id\}/messages}}

\paragraph{Amac}
AI'ya yeni mesaj gonderir ve cevap alir.

\paragraph{Request body}
\begin{lstlisting}
{
  "content": "THYAO icin 1 saatlik RSI tabanli bir strateji cikar",
  "context": {
    "user_id": "demo",
    "ticker": "THYAO",
    "timeframe": "1h",
    "indicator_id": "rsi",
    "selected_symbols": ["THYAO"],
    "auto_save_drafts": true
  }
}
\end{lstlisting}

\paragraph{Context alanlari}
\begin{itemize}[leftmargin=1.5em]
    \item \texttt{user\_id}
    \item \texttt{ticker}
    \item \texttt{timeframe}
    \item \texttt{indicator\_id}
    \item \texttt{active\_blueprint}
    \item \texttt{selected\_symbols}
    \item \texttt{auto\_save\_drafts}
\end{itemize}

\paragraph{Beklenen cevap}
Kesin shape workflow'a gore degisebilir, ancak tipik olarak su tur alanlar doner:

\begin{itemize}[leftmargin=1.5em]
    \item AI tarafinin yazdigi cevap metni
    \item onerilen aksiyonlar
    \item kullanilan tool sonuclari
    \item olusmus strategy/rule/indicator draft'i
    \item gerekiyorsa arka planda baslatilan backtest run id'si
\end{itemize}

\paragraph{Model davranisi}
\begin{itemize}[leftmargin=1.5em]
    \item \texttt{GEMINI\_API\_KEY} varsa Gemini istemcisi denenir
    \item aksi halde heuristic fallback devreye girebilir
    \item \texttt{GOOGLE\_CLOUD\_PROJECT} verilirse Vertex tabanli client denenebilir
\end{itemize}

\subsection{\texttt{GET /v1/ai/assets}}

\paragraph{Amac}
Belirli kullanicinin kayitli AI asset'lerini listeler.

\paragraph{Query parametresi}
\begin{itemize}[leftmargin=1.5em]
    \item \texttt{userId} (varsayilan: \texttt{demo})
\end{itemize}

\paragraph{Cevap}
\begin{lstlisting}
{
  "strategies": [],
  "rules": [],
  "indicators": []
}
\end{lstlisting}

\paragraph{Onemli not}
Cloud Run uzerinde varsayilan asset store yolu \texttt{/tmp/evalon\_ai\_assets.json} oldugundan, bu veri kalici bir veritabani yerine gecici instance dosyasi uzerinde tutulur.

\subsection{\texttt{POST /v1/ai/strategies}}
\subsection*{\texttt{POST /v1/ai/rules}}
\subsection*{\texttt{POST /v1/ai/indicators}}

\paragraph{Amac}
Kullaniciya ait AI uretimi draft asset kaydeder.

\paragraph{Ortak request body}
\begin{lstlisting}
{
  "userId": "demo",
  "title": "Momentum Blend",
  "description": "Draft indicator",
  "prompt": "RSI ile EMAyi birlestir",
  "spec": {
    "notes": ["ilk taslak"]
  }
}
\end{lstlisting}

\paragraph{Fark}
Path'e gore asset \texttt{strategy}, \texttt{rule} veya \texttt{indicator} olarak kaydedilir.

\section{Kullanim Senaryolari}

\subsection{Senaryo 1: Canli fiyat cekmek}

\begin{enumerate}[leftmargin=1.5em]
    \item \texttt{/v1/prices} cagir
    \item \texttt{ticker}, \texttt{timeframe}, gerekiyorsa \texttt{start/end/limit} ver
    \item cevapta \texttt{data[]} listesini chart'ta kullan
\end{enumerate}

\subsection{Senaryo 2: Indikator overlay cikarmak}

\begin{enumerate}[leftmargin=1.5em]
    \item once \texttt{/v1/indicators/catalog} ile desteklenen indikatorleri cek
    \item sonra \texttt{/v1/indicators} ile hedef ticker/timeframe icin seri hesaplat
    \item donen \texttt{indicators[]} verisini chart overlay veya panel olarak goster
\end{enumerate}

\subsection{Senaryo 3: Asenkron backtest}

\begin{enumerate}[leftmargin=1.5em]
    \item katalog endpointlerinden rule ve preset bilgilerini cek
    \item blueprint olustur
    \item \texttt{POST /v1/backtests/start} ile kuyruga al
    \item dondurulen \texttt{runId} ile status polling yap
    \item is bittiginde \texttt{/status}, \texttt{/events}, \texttt{/portfolio-curve} cagir
\end{enumerate}

\subsection{Senaryo 4: AI ile strateji taslagi cikarip kaydetmek}

\begin{enumerate}[leftmargin=1.5em]
    \item \texttt{POST /v1/ai/sessions} ile session ac
    \item \texttt{POST /v1/ai/sessions/\{id\}/messages} ile istegi gonder
    \item gerekirse tool katalogu ve backtest sonucuyla AI cevabini zenginlestir
    \item sonuc draft ise \texttt{/v1/ai/strategies}, \texttt{/v1/ai/rules} veya \texttt{/v1/ai/indicators} ile kaydet
\end{enumerate}

\section{Sinirlar ve Dikkat Edilmesi Gerekenler}

\begin{itemize}[leftmargin=1.5em]
    \item API su an auth gerektirmez; public kullanimda dikkat gerekir.
    \item Backtest run store memory tabanlidir; durable degildir.
    \item AI session store memory tabanlidir; durable degildir.
    \item AI asset store Cloud Run uzerinde \texttt{/tmp} tabanlidir; kalici veritabani yerine gecici depolama kullanir.
    \item \texttt{/v1/backtests/start} thread tabanlidir; yatay olceklenmis ortamlarda coordination mekanizmasi yoktur.
    \item \texttt{ticker=TEST} yalnızca test ve smoke test amacli ozel moddur.
    \item Oracle wallet veya sifre eksikse fiyat endpointleri \texttt{500} ile fail eder.
    \item Indikator hesaplama TA-Lib gerektirir.
\end{itemize}

\section{Web ve Mobile Entegrasyonu}

\subsection{Web}

Frontend tarafinda API base URL env uzerinden kullanilmalidir:

\begin{lstlisting}[language=bash]
NEXT_PUBLIC_EVALON_API_URL=https://evalon-backtest-api-474112640179.europe-west1.run.app
\end{lstlisting}

Mevcut Next.js katmaninda \texttt{/api/prices}, \texttt{/api/prices/batch} ve \texttt{/api/markets/list} route'lari bu URL'e proxy veya fetch yapacak sekilde kullanilir.

\subsection{Mobile}

Mobile tarafta build-time host verilmelidir:

\begin{lstlisting}[language=bash]
EVALON_API_BASE_URL=evalon-backtest-api-474112640179.europe-west1.run.app
\end{lstlisting}

Mobile tarafinda sadece host verilir; \texttt{https://} prefix'i kod tarafinda eklenir.

\section{Hizli Baslangic}

\subsection{Lokal smoke test}

\begin{lstlisting}[language=bash]
curl http://127.0.0.1:8000/health
curl "http://127.0.0.1:8000/v1/prices?ticker=TEST&timeframe=1m&limit=2"
curl "http://127.0.0.1:8000/v1/indicators/catalog"
\end{lstlisting}

\subsection{Cloud Run smoke test}

\begin{lstlisting}[language=bash]
bash scripts/smoke_test_cloud_run.sh \
  https://evalon-backtest-api-474112640179.europe-west1.run.app
\end{lstlisting}

\subsection{OpenAPI ve Swagger}

Servis kendi kendini belgeleyen iki onemli endpoint de sunar:

\begin{itemize}[leftmargin=1.5em]
    \item \texttt{/docs}
    \item \texttt{/openapi.json}
\end{itemize}

Bu nedenle client gelistirirken en guncel saha isimleri icin Swagger/OpenAPI ciktisi de kontrol edilmelidir.

\section{Sonuc}

Bu servis tek basina su is akisini kapsar:

\begin{itemize}[leftmargin=1.5em]
    \item BIST mum verisini cek
    \item indikator hesapla
    \item kural ve preset kataloglariyla strateji blueprint olustur
    \item senkron veya asenkron backtest kos
    \item event ve portfolio curve ciktilarini al
    \item AI tarafinda strateji/rule/indikator taslagi olustur ve kaydet
\end{itemize}

Uretim ortaminda bir sonraki mantikli adim durable storage, auth ve observability katmanlarini eklemektir. Ancak mevcut haliyle servis, chart, screener, backtest ve AI prototipleme ihtiyaclarini tek API uzerinden karsilayabilecek durumda tasarlanmistir.

\end{document}
